\section{Introduction}
\IEEEPARstart{D}{eep} Neural Networks (DNNs) have established foundational breakthroughs across core computer vision tasks, ranging from large-scale visual recognition \cite{8, 19} to dense spatial reasoning and object detection \cite{10, 37}. This technological trajectory was initially catalyzed by Convolutional Neural Networks (CNNs) \cite{15, 18}, which embedded translation equivariance and localized hierarchies into learnable architectures. To overcome gradient degradation in deeper layers, residual connections \cite{10} and compound scaling principles \cite{30} subsequently enabled networks to scale in depth and width. In recent years, the visual landscape has experienced an additional paradigm shift toward Vision Transformers (ViTs) \cite{6, 33}, masked image modeling \cite{34}, and self-supervised foundation backbones \cite{1}. While relaxing convolutional inductive biases enables dynamic global contextual reasoning, modern progress has predominantly reinforced an unsustainable ``scale-driven'' paradigm \cite{32}: marginal benchmark gains are increasingly acquired by expanding architectures to tens or hundreds of millions---and increasingly billions---of parameters over colossal compute clusters and proprietary web-scale datasets \cite{8, 32} (as analyzed across dominant paradigms in Section~\ref{sec:related_work}).

In parallel with these computational trends, deep learning has profoundly influenced clinical radiology, demonstrating diagnostic efficacy competitive with subspecialized clinicians across diverse imaging modalities \cite{7, 28}. Within neuro-oncology, non-invasive Magnetic Resonance Imaging (MRI) serves as the definitive clinical modality for the screening, differential diagnosis, and histologic grading of intracranial neoplasms \cite{2, 26, 38}. Clinical decision-making requires precise differentiation among distinct pathological entities---primarily diffuse infiltrative gliomas, extra-axial dural meningiomas, and sellar pituitary adenomas---as each exhibits vastly disparate treatment pathways ranging from urgent surgical resection and chemoradiotherapy to conservative observation \cite{26, 40}. Radiologically, this differential triage demands meticulous evaluation of subtle cytoarchitectural hallmarks: discerning irregular invasive margins from well-circumscribed borders, characterizing microvascular permeability and blood-brain barrier disruption, and distinguishing non-enhancing infiltrative edema from central necrotic cores across multi-parametric acquisitions (including T1-weighted, contrast-enhanced T1ce, T2-weighted, and FLAIR sequences) \cite{2, 28, 38}. However, manual multi-planar slice inspection across axial, coronal, and sagittal volumes imposes immense cognitive burdens on neuroradiologists, creating clinical vulnerability to intra- and inter-observer diagnostic variability and triage delays \cite{7, 40}. To automate the detection of these complex radiological phenotypes, prevailing medical imaging pipelines have routinely transposed off-the-shelf, overparameterized computer vision backbones into clinical workflows \cite{7, 39, 40}, inadvertently importing their extreme memory and computational footprints into hospital environments.

Despite these empirical successes, a critical translational divide persists between benchtop benchmark accuracy and clinical, point-of-care utility \cite{21, 25}. In healthcare environments, real-world deployment rarely occurs in centralized, high-throughput cloud data centers. Routine transmission of high-resolution volumetric neuroimaging to external cloud infrastructures encounters severe regulatory barriers, notably patient confidentiality mandates and data protection governance such as HIPAA in the United States and GDPR in the European Union \cite{21}. Furthermore, cloud dependency introduces vulnerability to variable network bandwidth, transmission latency, and service outages---risks that are unacceptable during acute clinical scenarios, emergency neuro-triage, or intra-operative surgical navigation \cite{21, 25}. Consequently, bringing intelligence directly to the clinical edge---operating locally on hospital workstations, embedded scanner consoles, and mobile intra-operative carts---has emerged as an urgent imperative \cite{9, 20, 21}.

However, edge deployment imposes uncompromising physical constraints: strict inference latency windows, severely constrained Random Access Memory (RAM), low thermal design power (TDP), and limited on-chip arithmetic hardware \cite{9, 25}. To accommodate mobile and decentralized environments, the vision community introduced lightweight architectural designs \cite{22} (detailed in Section~\ref{sec:related_work}), including depthwise separable convolutions \cite{11, 27}. Yet, the overwhelming majority of these ``efficient'' models predominantly inhabit a mobile-scale budget of 2 to 20 million parameters. In contrast, the sub-1M parameter regime ($<1,000,000$ parameters) remains fundamentally underexplored in medical imaging, particularly under the practical constraints of small-to-medium institutional datasets \cite{3, 16}. Crucially, as observed in recent foundational studies on compact architectures \cite{16}, naively downscaling standard architectures via uniform width or depth multipliers leads to erratic representational capacity, gradient instability, and disproportionate performance collapse \cite{16, 19}.

This investigation is motivated by the absence of systematic benchmarks, architectural insights, and deployment guidelines for deep learning models operating under a strict sub-1M parameter ceiling in neuro-oncological imaging. We posit that as the parameter budget is aggressively constrained, architectural topology---specifically how parameters are structurally allocated to capture multi-scale context without parameter redundancy---matters far more than sheer parameter volume \cite{16}. While multi-path atrous (dilated) depthwise convolutions have emerged as a premier mechanism for expanding effective receptive fields without parameter inflation \cite{5, 16}, their behavior, domain robustness, and failure mechanics under clinical distribution shifts have never been systematically characterized.

To address these foundational questions, this work presents a systematic empirical evaluation of deep learning models operating under a strict sub-1M parameter budget across representative neuro-oncological imaging tasks. Our empirical findings uncover two fundamental phenomena:
\begin{enumerate}
\item \textbf{In-Domain Saturation}: On canonical diagnostic benchmarks, architecturally optimized compact models match the diagnostic capability of overparameterized backbones an order of magnitude larger, while offering substantially faster inference speeds on unaccelerated edge processors. This demonstrates that for constrained clinical cohorts, high parameter capacity often yields diminishing returns.
\item \textbf{Cross-Dataset Multi-Scale Vulnerability}: When evaluated under out-of-domain clinical shifts, this saturation breaks. We find that compact models relying on parallel atrous convolutions become brittle under input scale perturbations, with predictive uncertainty concentrating along ambiguous lesion boundaries.
\end{enumerate}

In resource-constrained healthcare environments, resorting to larger models or heavy ensemble backbones to mitigate domain fragility is practically impermissible due to strict latency and memory limits \cite{9, 16}. Instead, we show that this residual accuracy gap can be recovered far more efficiently at inference time by leveraging the model's own architectural inductive biases. To this end, we introduce \textbf{Multi-Scale Consistency Gated Test-Time Augmentation (MCG-TTA)}. Rather than performing costly, unconditional multi-view inference across every image \cite{4, 35}, MCG-TTA monitors prediction uncertainty and selectively triggers targeted scale augmentations that align with the network's internal receptive field structure only when scale disagreement is detected, requiring zero additional learned parameters.

Extensive evaluations confirm that uncertainty-aware test-time calibration stabilizes compact representations under distribution shifts while preserving patient-level diagnostic reliability. These results demonstrate that architecture-informed inference calibration offers a practical, hardware-friendly alternative to model scaling for clinical edge artificial intelligence.

In summary, the principal contributions of this work are:
\begin{itemize}
\item \textbf{Systematic Sub-1M Benchmark and Edge Profiling}: We provide a comprehensive evaluation of sub-1M parameter vision architectures on neuro-oncological imaging, analyzing trade-offs across model capacity, real-device edge CPU latency, throughput, and cross-cohort generalization against mobile and overparameterized baselines, complemented by patient-level clinical sensitivity analysis.
\item \textbf{Characterization of Saturation and Domain Vulnerabilities}: We identify the saturation behavior of compact networks on in-domain medical diagnosis, and formally uncover the scale-sensitivity mechanics of compact atrous convolutions under clinical distribution shifts.
\item \textbf{Architecture-Informed Test-Time Inference}: We introduce MCG-TTA, a parameter-free, uncertainty-gated inference framework that dynamically stabilizes scale-sensitive compact representations at test time with minimal latency overhead.
\end{itemize}
